\documentclass[12pt,letterpaper]{article}

% ---------- packages ----------
\usepackage[T1]{fontenc}
\usepackage[utf8]{inputenc}
\usepackage{mathptmx}                 % Times-like text and math
\usepackage[margin=1in]{geometry}
\usepackage{amsmath,amssymb}
\usepackage{booktabs}
\usepackage{graphicx}
\usepackage{caption}
\usepackage{setspace}
\usepackage[round,authoryear]{natbib}
\usepackage{hyperref}
\hypersetup{colorlinks=true,linkcolor=black,citecolor=black,urlcolor=blue}
\usepackage{microtype}

\setlength{\parskip}{0pt}
\onehalfspacing

% ---------- title ----------
\title{\bfseries Admissibility Before Performance:\\
A Performance-Blind Structural Gate for\\
Regime-Adaptive Hierarchical Risk Parity}
\author{Om Patel\\[2pt]
\small Stevens Institute of Technology, Hoboken, NJ\\
\small \texttt{opatel4@stevens.edu}}
\date{\today}

\begin{document}
\maketitle

\begin{abstract}
\noindent Adaptive portfolio construction methods are conventionally evaluated by
out-of-sample performance. This creates a well-documented inferential problem: the
specification of the adaptive mechanism---thresholds, smoothing, event
definitions---is typically fixed after the researcher has observed how those
choices affect returns. We propose and apply an alternative protocol in which an
adaptive mechanism must first demonstrate \emph{structural informativeness} against
a prospectively frozen, performance-blind admissibility gate, and only then becomes
eligible for performance evaluation.

We instantiate this protocol for Regime-Adaptive Clustered Hierarchical Risk Parity
(RAC-HRP), in which re-clustering of the HRP dendrogram is triggered by changes in
the absorption ratio \citep{kritzman2011}. The gate evaluates four pre-registered
trigger thresholds against five hard structural criteria, none of which references
any return quantity. On a CRSP-native large-cap universe over a development region
of 2003--2022, \textbf{no threshold is admissible}. Triggered rebalances exhibit
greater clustering change than non-triggered rebalances at every threshold, and the
effect grows with selectivity, but no candidate attains statistical significance
after pre-registered multiplicity control (Holm-adjusted $p \in [0.221, 0.556]$).
Because the gate is conjunctive and performance-blind, the experiment terminates
without a portfolio ever being constructed and without the holdout sample being
opened.

A pre-registered post-mortem then asks whether the trigger's unusually bursty
timing is an artefact of the estimation pipeline. Across 1{,}500 replications of the
frozen machinery on regime-free inputs, real-data excess burstiness exceeds the
$97.5$th percentile of both an uncorrelated and a static-covariance null at all four
thresholds, so the tested regime-free architectural explanations are not supported.

Separately, a frozen estimator sensitivity sweep decomposes two components of
the modified HRP framework of \citet{molyboga2020}: within the development region, exponentially
weighted constant-correlation covariance improves every covariance-dependent
strategy examined, whereas the equal-volatility allocation modification remains
below static HRP under every covariance estimator.

We argue the resulting negative result is more informative than a comparable
performance-based null, precisely because no performance quantity was available at
the decision point. All code, frozen specifications, decision records, and a
hash-pinned reproducibility manifest are released.

\vspace{6pt}
\noindent\textbf{Keywords:} hierarchical risk parity; absorption ratio; regime
detection; pre-registration; multiple testing; reproducible research

\vspace{4pt}
\noindent\textbf{JEL classification:} G11, G17, C58, C52
\end{abstract}

\section{Introduction}

Hierarchical Risk Parity \citep{lopezdeprado2016} allocates capital using a
correlation-derived dendrogram rather than a matrix inversion, which makes it robust
in the high-dimensional regime where sample covariance is poorly conditioned. A
natural extension is to make the dendrogram \emph{adaptive}: rather than
re-clustering on a fixed calendar, re-cluster when the market's correlation
structure actually changes. The absorption ratio---the share of total variance
captured by the leading eigenvectors---is an appealing trigger because it is a
direct, model-free measure of covariance concentration and is known to move around
episodes of systemic fragility \citep{kritzman2011}.

This paper is about what happened when we tried to test that idea honestly.

\subsection{The inferential problem}

An adaptive mechanism of this kind is a \emph{specification family}, not a single
specification. The researcher must choose a smoothing window, a scale for
normalising changes, a firing threshold, and a re-clustering rule. Each choice
changes both the number of trigger events and, downstream, the portfolio's returns.
If the choices are made while performance is observable---even informally, even in
good faith---the resulting out-of-sample statistic is conditioned on the outcome it
purports to test.

The magnitude of this problem is not hypothetical. \citet{nikolopoulos2026} shows
that adaptive specification search generates statistically significant backtests
even under martingale-difference nulls, and proposes falsifying complete predictive
workflows against synthetic reference classes---including zero-predictability
environments---rather than auditing individual results. \citet{sheppert2026}
documents the corresponding degradation in reported strategy performance under
realistic search intensity, and \citet{bailey2014} give the general account of how
backtest overfitting propagates into out-of-sample claims. The concern is structural
rather than a matter of researcher discipline: \citet{demiguel2009} found that none
of fourteen optimisation methods they tested consistently beat naive $1/N$
out-of-sample, a result frequently attributed to estimation error but equally
consistent with specification search in the methods being compared.

\subsection{Admissibility before performance}

Our response is to separate two questions that are usually answered together:
\begin{enumerate}
\item \textbf{Is the adaptive mechanism structurally informative?} Do its trigger
      events identify moments when the dependence structure actually changes?
\item \textbf{Does acting on it improve performance?}
\end{enumerate}

We argue (1) should be settled first, using statistics that do not reference
returns, and that failure at (1) should terminate the investigation. A mechanism
that cannot demonstrate it is detecting the thing it claims to detect has no
principled basis for a performance claim, and a performance evaluation conducted
anyway will be uninterpretable---a positive result cannot be distinguished from
specification search, and a negative result cannot be distinguished from bad luck.

Concretely, we pre-register a \textbf{conjunctive structural gate} with five hard
criteria evaluated on a development region, freeze and hash it before any Phase~2
gate code is written, and implement it so that the selection rule is applied
automatically in code. The freeze is evidenced by the document hash recorded in the
released manifest together with the commit timestamps of the implementing code, not
by a countersignature: this is a single-author study and no independent party
authorised the specification (\S\ref{sec:no-independent-authorisation}). The Phase~2
calibration implementation never calls the return-performance path or risk-free
series and cannot produce a Sharpe ratio. Phase~0.5 did compute Sharpe ratios as part
of the falsification audit: environments A, B and D used null return environments,
while environment C retained development-region returns but randomised trigger
timing. Performance evaluation is reachable only through the gate.

\subsection{Contributions}

\paragraph{A performance-blind admissibility protocol.} We give a complete,
executable specification of a structural gate for adaptive re-clustering
mechanisms---informativeness, event sufficiency, timing variation, calendar
separation, and cluster informativeness---with frozen inference, pre-registered
multiplicity control, and a stopping rule that forbids selecting a ``least bad''
candidate.

\paragraph{A clean negative result for an absorption-ratio trigger.} No threshold in
the pre-registered grid is admissible. The binding failure is cluster
informativeness: the clustering-change effect is positive in direction at every
threshold and increases with selectivity, but is statistically unresolved under the
frozen inference procedure. Because the gate is performance-blind, this conclusion is
not conditioned on any return outcome.

\paragraph{A mechanism post-mortem that rules out two architectural explanations.}
The trigger fires in unusually tight temporal bursts. We pre-register a
pipeline-level null in which regime-free returns are pushed through the identical
frozen machinery, and show across 1{,}500 replications that neither an uncorrelated
nor a static-covariance environment reproduces the observed burstiness at any
threshold.

\paragraph{A decomposition of modified HRP.} A frozen estimator sweep separates two
of \citeauthor{molyboga2020}'s \citeyearpar{molyboga2020} modifications, finding that
the covariance modification helps in this sample while the allocation modification
does not.

\paragraph{A hash-pinned reproducibility record}, including a machine-checkable
assertion that the holdout sample was never opened.

\section{Related Work}

\paragraph{Hierarchical allocation.} \citet{lopezdeprado2016} introduced HRP,
motivated by the instability of quadratic optimisers under estimation error; the
mechanism is avoidance of covariance inversion. Subsequent work extends the tree
construction \citep{avellaneda2019,pergher2026} or the risk-budgeting rule
\citep{molyboga2020}. Earlier work established clustering as a preprocessing step
for allocation more generally \citep{leon2017}. Empirical replications are mixed:
\citet{dekovic2025} report $1/N$ dominating HRP on risk-adjusted return over
2005--2023 on S\&P~500 constituents (Sharpe 0.868 versus 0.534 on the full index)
while confirming HRP's volatility reduction of roughly one percentage point. We note
that their universe is the constituent set as of March 2023---the authors state the
limitation explicitly---so their backtest is subject to survivorship bias. Our
universe is constructed point-in-time from CRSP membership, which removes that
channel.

The general pattern---HRP reduces risk but underweights the high-volatility names
that drive trending markets---motivates adaptivity.

\paragraph{Regime detection and the absorption ratio.} \citet{kritzman2011}
introduced the absorption ratio as a systemic-risk measure and documented that
increases precede market drawdowns. Regime-switching approaches
\citep{hamilton1989,kimnelson1998,angtimmermann2012} establish that volatility and
correlation regimes are persistent and economically costly to ignore.
\citet{akioyamen2020} apply PCA to 48 macroeconomic series and recover a two-regime
crisis/non-crisis split by $k$-means with the regime count inferred from silhouette
width rather than assumed; \citet{horvathissa2023} detect regimes online through a
path-signature two-sample test; \citet{zhang2025} and \citet{kriuk2026} build
regime-aware allocation systems on sectoral and correlation-network structure
respectively.

Our trigger sits in this family but is deliberately minimal: a two-sided threshold on
the standardised change in a smoothed absorption ratio. This is by design. The
question we ask is not whether a sophisticated regime detector can be built, but
whether a simple, transparent, fully pre-specifiable one clears a structural
admissibility standard---and a minimal trigger makes the gate's verdict
interpretable in a way an elaborate one would not.

\paragraph{Covariance estimation.} \citet{ledoitwolf2004a} established linear
shrinkage toward structured targets; \citet{ledoitwolf2004b} propose the
constant-correlation target we adopt, and \citet{ledoitwolf2022} extend to nonlinear
shrinkage. \citet{molyboga2020} applies exponentially weighted covariance with
constant-correlation shrinkage in a managed-futures context; we source the decay
parameter $\alpha = 0.996$ to \citet{pafka2004}.
\citet{bongiorno2025} pursue learned covariance cleaning. Our sensitivity analysis
compares sample, linear-shrinkage, nonlinear-shrinkage, and exponentially weighted
constant-correlation estimators.

\paragraph{Evaluation methodology.} \citet{lopezdeprado2018} formalises purged,
embargoed cross-validation for financial time series; \citet{bailey2014} quantify the
deflation of reported Sharpe ratios under multiple testing. Our contribution is
orthogonal and complementary to both: purging addresses leakage \emph{within} an
evaluation and deflation adjusts a statistic \emph{after} it is computed, whereas the
gate addresses what is permitted to reach evaluation at all. The distinction matters
because deflation requires knowing the number of trials, which is precisely what
informal specification search leaves unrecorded.

\paragraph{Positioning.} Pre-registration is established in clinical and,
increasingly, psychological research, and has been proposed for finance, but
published applications in empirical portfolio construction remain rare. To our
knowledge this is the first study in this literature to pre-register a conjunctive
structural gate whose selection rule is executed in code, to hash the specification before writing the implementing code, and to report the resulting negative outcome
without amendment.

\section{Data and Universe}

The panel is constructed from CRSP daily stock files. Ordinary common shares (share
codes 10, 11) on NYSE, AMEX and Nasdaq (exchange codes 1, 2, 3) are retained.
Delisting returns are spliced following \citet{shumway1997}, with an audit record of
every splice. Prices recorded as bid/ask midpoints (negative in CRSP) are taken in
absolute value for market-capitalisation computation.

The investable universe at each rebalance is the largest $N = 100$ eligible names by
lagged market capitalisation, subject to a minimum history screen and a maximum
missing-observation fraction, constructed point-in-time so that no future membership
information enters. Rebalancing is monthly (21 trading days).

\paragraph{Region definition.} The development region runs 2003-01-08 to 2022-12-30
(5{,}031 trading days, 240 rebalances). The holdout region begins 2023-01-03. The
holdout was \emph{never opened} (Section~\ref{sec:holdout}).

The covariance window is selected by a pre-registered deterministic rule from
candidates $\{504, 756, 1260\}$ as a function of the median realised universe size;
for $N = 100$ this yields $W = 504$.

\section{Methodology}

\subsection{The trigger}

Let $\hat{\Sigma}_t$ be the covariance estimate over the trailing $W = 504$ trading
days ending at rebalance $t$, with eigenvalues
$\lambda_1 \geq \cdots \geq \lambda_N$. The absorption ratio is
\begin{equation}
AR_t \;=\; \frac{\sum_{i=1}^{k} \lambda_{i,t}}{\sum_{i=1}^{N} \lambda_{i,t}},
\label{eq:ar}
\end{equation}
with $k$ fixed at the Marchenko--Pastur count of the first eligible rebalance and
held for the entire run ($k = 15$). Fixing $k$ ex ante avoids a subtle look-ahead
channel: a $k$ that adapts to each window makes $AR$ partly a function of how many
components happen to clear the noise threshold at that date.

The series is smoothed over five rebalances, differenced, and standardised by a
strictly trailing scale:
\begin{align}
AR^{s}_t &= \tfrac{1}{5}\textstyle\sum_{j=0}^{4} AR_{t-j}, \label{eq:smooth}\\
\Delta AR_t &= AR^{s}_t - AR^{s}_{t-1}, \label{eq:delta}\\
\hat{\sigma}_t &= \mathrm{sd}\!\left(\Delta AR_{t-12},\dots,\Delta AR_{t-1}\right),
\quad \text{ddof} = 1,\ \text{min.\ 6 obs.} \label{eq:sigma}
\end{align}
A trigger fires when
\begin{equation}
\left| \Delta AR_t \right| \;>\; \gamma \cdot \hat{\sigma}_t,
\label{eq:trigger}
\end{equation}
two-sided with zero location and strict inequality. Because $\hat{\sigma}$ requires
six observations, the first seven rebalances can never fire; the eligible set is
$|E| = 240 - 7 = 233$.

The candidate grid is $\gamma \in \{0.5, 1.0, 1.5, 2.0\}$, fixed ex ante.
\textbf{$\gamma$ was not tuned.}

\subsection{The structural gate}

A candidate is admissible only if it satisfies all five criteria in
Table~\ref{tab:gate}. \textbf{No criterion references any return quantity.}
Turnover is computed and reported as a diagnostic but cannot pass or fail a
candidate.

\begin{table}[htbp]
\centering
\caption{The five hard structural criteria. All are evaluated on the development
region; none references a return quantity.}
\label{tab:gate}
\small
\begin{tabular}{lll}
\toprule
Criterion & Statistic & Pass rule \\
\midrule
Informativeness & $f_\gamma = |T_\gamma| / |E|$ & $0.05 \leq f_\gamma \leq 0.40$ \\
Event sufficiency & events per development fold & $\geq 3$ in all four folds \\
Timing variation & CV of gaps; modal-gap share & CV $\geq 0.50$ \emph{and} modal $\leq 0.50$ \\
Separation & $J^{*}_\gamma$, max.\ Jaccard overlap & $\leq$ placebo 95th pct.\ at own $|T_\gamma|$ \\
Cluster informativeness & $D_{VI}$ (Eq.~\ref{eq:dvi}) & $D_{VI} > 0$ \emph{and} Holm $p < 0.05$ \\
\bottomrule
\end{tabular}
\end{table}

The \emph{separation} criterion deserves comment. A trigger that fires on a fixed
calendar is not regime-adaptive, whatever its label. We therefore compute
\begin{equation}
J^{*}_\gamma \;=\; \max_{q \in \{2,\dots,12\}} \; \max_{r \in \{0,\dots,q-1\}}
\frac{\left| T_\gamma \cap P(q,r) \right|}{\left| T_\gamma \cup P(q,r) \right|},
\label{eq:jstar}
\end{equation}
where $P(q,r)$ is the periodic schedule of period $q$ and phase $r$. Maximising over
both period \emph{and phase} is essential: a trigger firing every second rebalance on
odd-indexed dates scores $J^{*} = 0.252$ without phase adjustment (passing) and
$1.000$ with it (caught). The threshold is recomputed by Monte Carlo at each
candidate's own event count, because it depends strongly on that count---the 95th
percentile is $0.171$ at 12 events and $0.382$ at 112.

The \emph{cluster informativeness} criterion is the heart of the gate. It asks
whether triggered rebalances coincide with larger changes in the hierarchical
clustering than non-triggered ones, measured by variation of information ($VI$)
between consecutive clusterings on the intersection of consecutive universes:
\begin{equation}
D_{VI} \;=\; \mathrm{median}\!\left(VI_t \mid I_t = 1\right)
           - \mathrm{median}\!\left(VI_t \mid I_t = 0\right).
\label{eq:dvi}
\end{equation}
Critically, the counterfactual clustering is recomputed at \emph{every} eligible
rebalance, including those a live strategy would skip; without this, triggered and
non-triggered structural change are not comparable.

\subsection{Frozen inference}

All inference parameters were fixed before execution: placebo seed 20260817 with
100{,}000 draws; circular block bootstrap with 10{,}000 replicates and
Politis--White automatic block length \citep{politiswhite2004}, resampling the pair
$(VI_t, I_t)$ jointly in blocks to preserve alignment; one-sided
$H_0\!:\ D_{VI} \leq 0$ with replicates centred on the observed statistic; Holm
step-down across the four candidates at family-wise $\alpha = 0.05$. Holm rather than
Bonferroni because the trigger sets are highly dependent across $\gamma$.

\subsection{The stopping rule}

Among candidates passing \emph{every} hard criterion, select the value closest to the
inherited $\gamma = 1.0$, breaking ties toward the larger $\gamma$. \textbf{If none
passes, the investigation stops; the ``least bad'' candidate is not selected.} This
rule is applied automatically in code, not by the researcher.

\section{Results}

\subsection{Static baselines and estimator sensitivity}

Table~\ref{tab:baselines} reports development-region performance for five static
baselines under the pre-registered nonlinear-shrinkage estimator. The equal risk
contribution portfolio follows \citet{roncalli2013}; MHRP\_EV implements component~(ii)
of \citeauthor{molyboga2020}'s modification (equal-volatility rather than
inverse-variance budgeting) in isolation.

\begin{table}[htbp]
\centering
\caption{Static baselines, development region (2003--2022, 5{,}031 days).}
\label{tab:baselines}
\small
\begin{tabular}{lrrrrrr}
\toprule
Strategy & Ann.\ return & Ann.\ vol & Sharpe & Max DD & Calmar & Ann.\ turnover \\
\midrule
MinVar     & 0.092 & 0.154 & 0.576 & $-0.426$ & 0.215 & 1.502 \\
HRP\_static & 0.096 & 0.163 & 0.574 & $-0.461$ & 0.208 & 1.455 \\
ERC        & 0.096 & 0.173 & 0.549 & $-0.502$ & 0.191 & 0.656 \\
MHRP\_EV   & 0.093 & 0.175 & 0.534 & $-0.511$ & 0.183 & 1.779 \\
EW         & 0.094 & 0.191 & 0.508 & $-0.560$ & 0.168 & 0.640 \\
\bottomrule
\end{tabular}
\end{table}

An independent accounting reconciliation---recomputing the portfolio return series
from the stored weight path and realised asset returns, without reusing the engine's
arithmetic---matches the engine's gross returns to a median absolute difference of
$0.0$ for all five strategies.

Table~\ref{tab:estimators} and Figure~\ref{fig:estimators} report the estimator
sensitivity sweep. Equal weighting is invariant across estimators ($0.508$), as it
must be: it never consults the covariance matrix. This serves as a negative control
confirming that the observed differences are attributable to the estimator.

\begin{table}[htbp]
\centering
\caption{Estimator sensitivity (Sharpe). Diagnostic only; selects nothing.}
\label{tab:estimators}
\small
\begin{tabular}{lrrrr}
\toprule
Strategy & \texttt{ewma\_cc} & \texttt{lw\_linear} & \texttt{nls} & \texttt{sample} \\
\midrule
ERC        & 0.558 & 0.549 & 0.549 & 0.549 \\
EW         & 0.508 & 0.508 & 0.508 & 0.508 \\
HRP\_static & \textbf{0.608} & 0.591 & 0.574 & 0.594 \\
MHRP\_EV   & 0.568 & 0.554 & 0.534 & 0.557 \\
MinVar     & 0.588 & 0.551 & 0.576 & 0.556 \\
\bottomrule
\end{tabular}
\end{table}

Two findings follow, both bounded to the development region. First, exponentially
weighted constant-correlation covariance improves every covariance-dependent
strategy examined, most strongly HRP\_static ($0.574 \rightarrow 0.608$); it is the
largest estimator effect in the table. Second, the equal-volatility allocation
(MHRP\_EV) remains below static HRP under \emph{every} estimator. Since the ordering
is invariant to the estimator, the underperformance is attributable to the
allocation rule rather than to a covariance interaction. Taken together these
separate two components of \citeauthor{molyboga2020}'s construction: the covariance
modification transfers to this equity universe; the allocation modification does not.

We note that \texttt{ewma\_cc} outperforms the pre-registered \texttt{nls}
specification on four of five strategies. The pre-registered specification is
therefore not the ex-post best-performing estimator in this sample. \textbf{The
design is not retroactively changed to exploit this.}

\begin{figure}[htbp]
\centering
\includegraphics[width=0.75\textwidth]{fig4_estimator_sensitivity.pdf}
\caption{Estimator sensitivity across five static baselines. EW is invariant
because it never consults the covariance matrix, serving as a negative control.}
\label{fig:estimators}
\end{figure}

\subsection{The structural gate: no admissible threshold}

Figure~\ref{fig:ar} shows the absorption-ratio series over the development region
together with the standardised change and the four candidate thresholds;
Table~\ref{tab:gateresult} gives the gate result and Figure~\ref{fig:dvi}
summarises the cluster-informativeness outcome.

\begin{figure}[htbp]
\centering
\includegraphics[width=\textwidth]{fig1_absorption_ratio_and_triggers.pdf}
\caption{Absorption ratio (top) and standardised absorption-ratio change with the
four candidate thresholds (bottom), development region, $k = 15$ fixed, $W = 504$.}
\label{fig:ar}
\end{figure}

\begin{table}[htbp]
\centering
\caption{Frozen structural calibration gate. 233 eligible rebalances, $k = 15$.}
\label{tab:gateresult}
\small
\setlength{\tabcolsep}{4pt}
\begin{tabular}{crrrrrrrrrc}
\toprule
$\gamma$ & events & firing & min ev/fold & CV gap & modal gap & $J^{*}$ & $J$ thr & $D_{VI}$ & $p$ Holm & Adm. \\
\midrule
0.5 & 149 & 0.639 & 35 & 1.116 & 0.824 & 0.400 & 0.446 & $+0.033$ & 0.556 & No \\
1.0 & 111 & 0.476 & 26 & 1.458 & 0.791 & 0.326 & 0.382 & $+0.022$ & 0.556 & No \\
1.5 &  81 & 0.348 & 18 & 1.676 & 0.763 & 0.263 & 0.320 & $+0.080$ & 0.254 & No \\
2.0 &  58 & 0.249 & 11 & 1.823 & 0.791 & 0.208 & 0.259 & $+0.096$ & 0.221 & No \\
\bottomrule
\end{tabular}
\end{table}

No candidate is admissible. The pattern of failure is informative.

\paragraph{Separation passes everywhere.} $J^{*} \leq$ threshold at every $\gamma$.
The trigger is not a disguised fixed schedule---a substantive negative finding, since
that is the most common failure mode for signals of this kind.

\paragraph{Event sufficiency passes everywhere} (11--35 events per fold, all
$\geq 3$).

\paragraph{Informativeness fails only at $\gamma \in \{0.5, 1.0\}$}, where firing
rates of $0.639$ and $0.476$ exceed the $0.40$ ceiling. The inherited $\gamma = 1.0$
fires on nearly half of all eligible rebalances, which is difficult to reconcile with
a description of the events as unusual.

\paragraph{Timing variation fails everywhere, on modal-gap share alone.} Every CV
clears its floor comfortably ($1.116$--$1.823$); the failure is that modal-gap share
($0.763$--$0.824$) never approaches the $0.50$ ceiling. Section~\ref{sec:mechanism}
examines this.

\paragraph{Cluster informativeness fails everywhere.} $D_{VI}$ is positive at every
$\gamma$ and increases with selectivity ($+0.033 \rightarrow +0.096$), the direction
the hypothesis predicts. But no Holm-adjusted $p$ clears $0.05$ (best $0.221$ at
$\gamma = 2.0$). The defensible statement is that the effect is \textbf{positive in
direction but statistically unresolved}: a positive point estimate does not establish
that the effect is real.

\begin{figure}[htbp]
\centering
\includegraphics[width=\textwidth]{fig2_cluster_informativeness.pdf}
\caption{Cluster-informativeness result. Left: $D_{VI}$ point estimates with event
counts. Right: raw and Holm-adjusted $p$-values against the pre-registered
$\alpha = 0.05$.}
\label{fig:dvi}
\end{figure}

Even the strongest candidate, $\gamma = 2.0$---which passes informativeness and
separation---is blocked by two criteria. Because failure is conjunctive and because
cluster informativeness fails independently at every $\gamma$, the conclusion does
not rest on the contested timing criterion.

\medskip
\noindent\textbf{Per the frozen stopping rule, the investigation terminates here. No
portfolio was constructed and no performance quantity was computed.}

\subsection{Post-mortem: is the burstiness an artefact of the pipeline?}
\label{sec:mechanism}

The timing failure has a specific character: high CV \emph{and} high modal-gap share
together mean most inter-event gaps are one rebalance with occasional long quiescent
stretches---the trigger fires in bursts. Two mechanisms could produce this. Real
markets may cluster covariance-concentration episodes in time; or the estimation
pipeline may manufacture bursts from overlapping windows, smoothing, differencing,
and a rolling scale denominator.

A placement null (random reallocation of the same number of events across the same
grid) rejects the hypothesis that the burstiness is a mere consequence of event
density: observed modal-gap share exceeds the placement null at every $\gamma$. But a
placement null cannot separate the two mechanisms above, because it never exercises
the pipeline.

We therefore pre-registered a \textbf{pipeline-level null}, frozen and hashed before execution, with the decision rule fixed in advance and applied mechanically. Three
environments are pushed through the \emph{identical unmodified} frozen machinery,
preserving the real NaN mask, membership path, market caps, rebalance dates, $W$,
$k$, smoothing and $\hat{\sigma}$:
\begin{itemize}
\item \textbf{A}---vol-matched Gaussian, zero cross-correlation, no time dependence
      (floor control);
\item \textbf{S}---one static covariance $\Sigma_0$ fitted on development rows only,
      drawn i.i.d.\ through time: persistent cross-sectional structure, no regimes
      (adjudicating null);
\item \textbf{D}---two-state Markov process in which both volatility \emph{and}
      factor loadings scale, so covariance \emph{concentration} genuinely shifts
      (positive control).
\end{itemize}

Because modal-gap share depends on event count, comparisons use a density-adjusted
statistic
\begin{equation}
B_\gamma \;=\; M_\gamma - \mathrm{median}\!\left(M^{\mathrm{placement}} \mid n_\gamma\right),
\label{eq:bgamma}
\end{equation}
where the placement median is a deterministic cached lookup keyed on the event count
and shared with the real data. We run 500 replications per environment (1{,}500
total).

\begin{table}[htbp]
\centering
\caption{Pipeline-level null. Excess burstiness $B_\gamma$, 500 replications per
environment. Intervals are 2.5th--97.5th percentiles.}
\label{tab:mechanism}
\small
\begin{tabular}{crccc}
\toprule
$\gamma$ & real $B$ & A & S & D \\
\midrule
0.5 & 0.182 & 0.062--0.157 & 0.033--0.117 & 0.045--0.166 \\
1.0 & 0.318 & 0.122--0.272 & 0.100--0.237 & 0.140--0.306 \\
1.5 & 0.413 & 0.140--0.361 & 0.095--0.296 & 0.163--0.407 \\
2.0 & 0.526 & 0.107--0.392 & 0.000--0.295 & 0.105--0.440 \\
\bottomrule
\end{tabular}
\end{table}

\begin{figure}[htbp]
\centering
\includegraphics[width=\textwidth]{fig3_mechanism_null.pdf}
\caption{Real-data excess burstiness (horizontal line) against regime-free pipeline
nulls, 500 replications per environment, by $\gamma$.}
\label{fig:mechanism}
\end{figure}

At every $\gamma$ the real value exceeds the $97.5$th percentile of both A and S,
satisfying the pre-registered condition for the ``beyond regime-free mechanics''
outcome. The tested regime-free architectural explanations are therefore not
supported.

Three qualifications. First, D does not bracket the real data either, so the optional
consistency-with-nonstationarity statement is unavailable; but since D's states last
approximately 100 and 33 daily observations against $W = 504$, a 504-day window spans
roughly 3.8 complete cycles and materially averages the designed regimes away. D is a
\emph{high-frequency} positive control, and its non-overlap indicates that rapid
regime switching is heavily attenuated by the estimator rather than that regime
explanations are excluded. Second, the D margins are thin ($+0.006$ at
$\gamma = 1.5$) and should not carry independent weight. Third, S produced
\emph{less} excess burstiness than A at every $\gamma$---contrary to the hypothesis
that motivated S as the adjudicating null. We report this without explanation;
establishing one would require further work.

\subsection{Empirical size of the inference procedure}

Because the cluster-informativeness criterion is the binding constraint, its
calibration matters. We measured the empirical rejection rate under the null over
2{,}000 synthetic datasets at three replicate counts (Table~\ref{tab:size}).

\begin{table}[htbp]
\centering
\caption{Empirical size at nominal $\alpha = 0.05$, 2{,}000 null datasets.}
\label{tab:size}
\small
\begin{tabular}{rrrc}
\toprule
$B$ & rate & MC SE & 95\% CI \\
\midrule
600     & 0.0655 & 0.0055 & 0.0547--0.0763 \\
2{,}000  & 0.0660 & 0.0056 & 0.0551--0.0769 \\
10{,}000 & 0.0655 & 0.0055 & 0.0547--0.0763 \\
\bottomrule
\end{tabular}
\end{table}

The rate is stable across a seventeen-fold increase in replicates, so the departure
is a property of the procedure rather than a finite-$B$ artefact. The test is
\textbf{mildly anti-conservative}: empirical size $\approx 6.6\%$ against nominal
$5\%$. This strengthens rather than weakens the central result. The
cluster-informativeness test is biased \emph{toward} rejection, and no candidate
rejected.

\section{Discussion}

\subsection{What the negative result establishes}

Under a prospectively frozen, performance-blind admissibility standard, an
absorption-ratio change trigger does not qualify as an adaptive re-clustering
mechanism on this universe and development region. The trigger is not equivalent to
a fixed schedule, and its events are genuinely bursty relative to both random
placement and two regime-free pipeline nulls. But triggered rebalances are not shown
to identify moments of greater hierarchical restructuring.

The interesting asymmetry is that the trigger evidently responds to \emph{something}.
Its timing is far from random, and its burstiness survives two architectural
explanations. What it does not do is align with the quantity the adaptive hypothesis
requires. A signal can react strongly to real covariance events without those events
being moments when re-clustering changes the tree in a statistically reliable
way---for instance if $AR$ responds to concentration changes that leave the dendrogram
topology intact, or if it responds at a horizon offset from the restructuring.

\subsection{Why performance-blindness matters here}

Had this study measured performance first, the natural path would have been: pick
$\gamma = 1.0$ (inherited), observe results, and---on encountering a $47.6\%$ firing
rate---adjust the threshold. Each adjustment is individually defensible and
collectively fatal, in exactly the way \citet{nikolopoulos2026} quantifies.

The gate makes that path unavailable. The $0.40$ firing ceiling was fixed before the
firing rate was known and was \emph{not} revised when $\gamma = 1.0$ exceeded it. The
$\gamma$ grid was fixed before any event count was computed. The stopping rule was
executed in code. There is no point at which a performance number could have
influenced a specification choice, because no performance number existed.

\subsection{The Molyboga decomposition}

The estimator sweep produces an unanticipated positive finding.
\citet{molyboga2020} proposes three modifications to HRP; two are separately
implementable here. Within the development region, the covariance modification
improves every covariance-dependent strategy tested, whereas the allocation
modification remains below static HRP under every covariance estimator. This suggests
reported gains from ``modified HRP'' may be attributable to the covariance treatment
rather than the allocation rule---a decomposition obscured when the modifications are
evaluated jointly.

We emphasise this is in-sample within the development region. The holdout was not
opened, and no out-of-sample claim is made.

\section{Limitations}

\subsection{Statistical power}

The result is a failure to reject, not a demonstration of absence. With 233 eligible
rebalances and four folds, the design has limited power against small effects, and
$D_{VI}$ is positive at every threshold. A larger universe, a longer sample, or a
higher rebalance frequency could plausibly resolve the effect. The honest claim is
that this evidence is insufficient, not that the effect is zero.

\subsection{Horizon alignment between trigger and criterion}

The trigger statistic measures change in the absorption ratio over five rebalances
(Section~\ref{sec:trigger}) while the cluster-informativeness criterion measures change in the
clustering over one. The pre-registration did not require these horizons to agree, and
the gate was frozen without the mismatch being identified. The negative result
therefore establishes that the trigger is uninformative about one-rebalance
restructuring under the frozen procedure; it does not establish that the trigger is
uninformative about hierarchical structure generally. A subsequent pre-registration of
this protocol should specify the trigger and criterion horizons jointly.

\subsection{Single universe and single trigger family}

One universe (CRSP large-cap, $N = 100$), one region, one trigger construction. The
absorption-ratio trigger is one member of a large family; failure here does not
generalise to eigenspace-rotation, correlation-distance, or change-point triggers.

\subsection{The holdout was never opened}
\label{sec:holdout}

The single pre-registered holdout touch was not used. A machine-checkable assertion in
the released test suite verifies that the durable audit log records zero analysis
touches of the test region. This is a consequence of the design, not an omission:
performance evaluation was \emph{conditionally inaccessible}, and no adaptive
specification earned access.

\subsection{Reproducibility incidents}

Two defects were found and are documented in full. A figure in the frozen
pre-registration cited a Phase~0.5 re-clustering count as though it were a trigger
count (corrected by erratum; no conclusion affected). Separately, a commit deleted the
equal-volatility allocator while the driver still declared it, rendering an archived
result temporarily unreproducible; the code was restored from history and the archived
value reproduced exactly. Both are recorded in the released supplement, and regression
tests now cover both failure modes.

\subsection{No independent authorisation}
\label{sec:no-independent-authorisation}

This is a single-author study. The pre-registration was written, frozen, and applied by
the same person, and no independent party ruled on any specification choice or
amendment. The protocol's commitments are therefore evidenced by artefacts rather than
by review: specification documents are hashed and their hashes recorded in a released
manifest, the implementing code carries commit timestamps that postdate those hashes,
the selection rule is executed by a function rather than applied by hand, and the
calibration implementation cannot compute a performance quantity.

These devices constrain the specific failure they were built for -- selecting a
threshold after observing its effect on returns -- because the gate is structurally
incapable of producing a return quantity and the decision rule is executed in code.
They do not constrain judgement calls made between phases, where a self-imposed rule
can be reinterpreted by the person who wrote it. One such decision is reported in full
in Section~\ref{sec:phase05}, with both the original and revised outcomes, so that a reader may apply
either rule.

Independent authorisation would strengthen the protocol, and any subsequent application
of it should include a party who does not benefit from the outcome.

\section{Conclusion}

We proposed and applied a protocol in which an adaptive portfolio mechanism must
demonstrate structural informativeness against a frozen, performance-blind gate before
performance evaluation becomes available. Applied to an absorption-ratio
re-clustering trigger, the gate returned no admissible threshold, and the
investigation terminated without a portfolio being constructed and without the holdout
being opened.

The mechanism's timing is not random and is not reproduced by two regime-free pipeline
nulls, so it is detecting something. What it is not shown to detect is statistically
distinguishable change in the hierarchical dependence structure---the quantity the
adaptive hypothesis requires.

We suggest the value of this result lies less in the specific mechanism than in the
demonstration that a pre-registered structural standard can be specified precisely
enough to be executed automatically, and can return a negative answer that survives
adversarial review without being rewritten. The natural next question is not how to
make this trigger pass, but what $|\Delta AR|$ is actually tracking---an alignment
question best addressed by comparing absorption-ratio changes against eigenspace
rotation, correlation-matrix distance, and tree instability directly.

\section*{Acknowledgments}

The author thanks [names, if anyone read drafts or discussed the work] for helpful
discussion. This is a single-author study; no independent party reviewed or authorised
the methodological decisions described here, and the protocol's commitments rest on
hashed specifications and commit timestamps rather than on external review
(\S\ref{sec:no-independent-authorisation}).

The author used a large language model to assist with drafting, literature review,
auditing the project's own records, and identifying methodological defects. The author
reviewed and verified all output and is responsible for the content.

Any remaining errors are the author's own.

\section*{Data and Code Availability}

All code, frozen specifications, decision records, errata, and result artefacts are
available at \url{https://github.com/opatel4/rac-hrp}. A reproducibility manifest pins
SHA-256 hashes for 30 code modules, 6 result artefacts, 11 governing documents, and
the underlying data files, together with the software environment. CRSP data are
licensed and cannot be redistributed; the manifest records checksums so a licensed
replicator can verify identity.

\bibliographystyle{plainnat}
\bibliography{references}

\end{document}
